% =====================================================================
%  Section: Method  (IJCAI 2027 revision: design-rationale centric)
% =====================================================================
\section{Method}
\label{sec:method}

\paragraph{Design philosophy.}
A single painting-animation system must simultaneously satisfy four
\emph{conflicting} objectives: perceptual grouping of artistically rendered
content, geometric reconstruction of objects that have no real 3D extent,
user-specifiable composition, and temporally coherent stylized rendering.
Jointly optimizing a single network for all four is ill-posed because the
gradients of style fidelity and 3D structure pull in opposite directions.
Our central design principle is therefore \emph{modular separation of
concerns with explicit handshakes}: each objective is addressed by an
independently optimized specialist, and the four specialists communicate
only through well-defined, interpretable intermediate
representations---masks with depth-ordered layers, a metric-scaled
object-level 3D library, a composited 3D scene with disocclusion-filled
background, and finally a multi-view keyframe stream with paired
geometry and edge cues. This decomposition converts a conflicting joint
problem into four well-posed sub-problems, while the handshakes carry
exactly the information each downstream stage needs and no more,
preventing error compounding.

\paragraph{Problem statement.}
Given a painting scan $P \in \mathbb{R}^{H\times W \times 3}$ and a set of
hand-drawn line-art sketches $L = \{\ell_1, \dots, \ell_K\}$ specifying
target object positions, Canvas3D produces an animated video $V$ that
(i) preserves the artistic style of $P$, (ii) supports viewpoint
exploration of a 3D-composited scene, and (iii) places each foreground
object at its sketched location with geometrically grounded, temporally
coherent motion. The framework comprises four modules
(Figure~\ref{fig:overview}): \textbf{PaintParse} (perceptual grouping and
layer decomposition), \textbf{MonoSplat} (per-object 3D lifting and
library assembly), \textbf{SketchComp} (line-art-guided scene composition
with depth-aware inpainting), and \textbf{ArtiFlow} (geometry-grounded
video synthesis).

\paragraph{Notation.}
$M_i \in \{0,1\}^{H\times W}$ denotes a binary mask for the $i$-th
foreground object, $M_\text{bg}$ the residual background mask,
$D_\text{mono}$ the monocular depth map, $\mathcal{G}_i$ the 3D
representation of object $i$, $\mathcal{S}$ the composited scene,
$\Omega_\text{gap}$ the disocclusion region, and $\{I_t, D_t, E_t\}$ the
multi-view RGB, depth, and edge keyframes.

\begin{figure}[t]
    \centering
    \resizebox{\linewidth}{!}{%
    \begin{tikzpicture}[
        font=\scriptsize,
        stage/.style={draw=black!60, rounded corners=4pt, thick,
            minimum width=2.7cm, minimum height=1.6cm, align=center},
        inp/.style={draw=black!50, fill=yellow!10, rounded corners=3pt,
            minimum width=1.8cm, minimum height=0.7cm, align=center,
            font=\scriptsize},
        outp/.style={draw=black!50, fill=gray!8, rounded corners=3pt,
            minimum width=2.2cm, minimum height=0.7cm, align=center,
            font=\scriptsize},
        ar/.style={->, >=Stealth, thick, draw=black!55},
        feedback/.style={->, >=Stealth, thick, draw=red!65, dashed}
    ]
        % Inputs (top row)
        \node[inp] (pin) {Painting $P$};
        \node[inp, right=0.7cm of pin] (lin) {Line-art $L$};

        % Stage 1: PaintParse
        \node[stage, fill=blue!8, below=0.45cm of pin] (pp)
            {\textbf{PaintParse}\\[2pt]
             \footnotesize perceptual grouping\\
             \footnotesize + layer ordering};
        \node[outp, below=0.18cm of pp] (ppo)
            {$\{M_i\},\;M_\text{bg},\;\{d_i\}$};

        % Stage 2: MonoSplat
        \node[stage, fill=green!10, below=0.45cm of ppo] (ms)
            {\textbf{MonoSplat}\\[2pt]
             \footnotesize per-object 3D lifting\\
             \footnotesize + metric rescaling};
        \node[outp, below=0.18cm of ms] (mso)
            {3D library $\{\mathcal{G}_i\}$};

        % Stage 3: SketchComp
        \node[stage, fill=orange!12, below=0.45cm of mso, xshift=1.4cm] (sc)
            {\textbf{SketchComp}\\[2pt]
             \footnotesize contour registration\\
             \footnotesize + depth inpainting};
        \node[outp, below=0.18cm of sc] (sco)
            {Scene $\mathcal{S}^*,\,\mathcal{D}^*$};

        % Stage 4: ArtiFlow
        \node[stage, fill=red!10, below=0.45cm of sco] (af)
            {\textbf{ArtiFlow}\\[2pt]
             \footnotesize multi-view rendering\\
             \footnotesize + dual-conditioned diffusion};
        \node[outp, below=0.18cm of af] (afo)
            {Video $V$};

        % Arrows
        \draw[ar] (pin) -- (pp);
        \draw[ar] (pp) -- (ppo);
        \draw[ar] (ppo) -- (ms);
        \draw[ar] (ms) -- (mso);
        \draw[ar] (mso) -- (sc);
        \draw[ar] (lin) |- ++(0,-3.4) -| (sc);
        \draw[ar] (sc) -- (sco);
        \draw[ar] (sco) -- (af);
        \draw[ar] (af) -- (afo);

        % Feedback loop
        \draw[feedback] (af.west) -- ++(-1.4,0)
            node[midway, above, font=\tiny, text=red!75] {re-composition}
            |- (sc.west);

        % Stage labels on the right
        \node[right=0.1cm of pp, font=\tiny, text=blue!60!black, align=left]
            {concern:\\perception};
        \node[right=0.1cm of ms, font=\tiny, text=green!60!black, align=left]
            {concern:\\geometry};
        \node[right=0.1cm of sc, font=\tiny, text=orange!70!black, align=left]
            {concern:\\composition};
        \node[right=0.1cm of af, font=\tiny, text=red!65!black, align=left]
            {concern:\\temporal};
    \end{tikzpicture}}
    \caption{\textbf{Canvas3D overview.} A modular pipeline that factors
    painting animation into four independently optimized specialists:
    \emph{PaintParse} (perceptual grouping and depth-ordered layering),
    \emph{MonoSplat} (per-object 3D lifting into a reusable library),
    \emph{SketchComp} (line-art-guided scene composition with
    depth-aware inpainting), and \emph{ArtiFlow} (geometry-grounded
    video synthesis with dual control signals). A dashed feedback loop
    re-composes the scene when per-object motion verification fails.}
    \label{fig:overview}
\end{figure}

% =====================================================================
\subsection{PaintParse: Art-Domain Perceptual Grouping and Layer Ordering}
\label{sec:method:paintparse}

\paragraph{Design rationale.}
The downstream 3D and composition stages require two distinct outputs
that no single head of an off-the-shelf detector provides:
\emph{(i)} foreground masks whose boundaries respect artistic intent
rather than photographic edges, and \emph{(ii)} a depth-consistent
layer ordering that disambiguates which object lies in front of which.
PaintParse therefore couples a domain-adapted instance segmentation
backbone with a parallel monocular depth estimator and fuses their
outputs through a percentile-based layer assignment. The key design
choice is \emph{decoupling semantic grouping from depth ordering}: the
segmentation network is responsible only for ``what is an object,''
while the depth estimator is responsible only for ``how far away it
is,'' and the layer-assignment logic simply reads both off. This
avoids asking either network to solve a problem it was not trained
for.

\paragraph{Instance segmentation with art-domain adaptation.}
We adopt a one-stage instance segmentation backbone pre-trained on
natural-image data. Direct transfer to paintings yields poor mask
quality because brushstroke textures produce spurious edge responses
and depth-ambiguous compositions confound object-background separation.
We mitigate this through a two-stage adaptation protocol: a
\emph{low-rank adaptation} of the detection backbone to inject
art-domain statistics without catastrophic forgetting, followed by
\emph{full fine-tuning} of the mask-prediction head to align boundary
prediction with painterly edges. Given input $P$, the adapted model
outputs $K$ binary masks $\{M_i\}_{i=1}^{K}$. The residual background
mask is computed by pixel-wise subtraction:
\begin{equation}
    M_\text{bg}(p) \;=\; 1 \,-\, \max_{i\in\{1,\dots,K\}} M_i(p),
    \quad \forall p \in [1,H]\!\times\![1,W].
    \label{eq:bg_mask}
\end{equation}

\paragraph{Monocular depth and layer assignment.}
In parallel, a monocular depth estimator produces a relative depth map
$D_\text{mono}$. For each detected object $i$, we summarize its depth
by the 10th and 90th percentiles of $\{D_\text{mono}(p)\mid p\in\Omega_i\}$
where $\Omega_i = \{p \mid M_i(p){=}1\}$, yielding a robust depth range
$[d_i^\text{near}, d_i^\text{far}]$ that is insensitive to outlier
pixels at mask boundaries. The layer index of each object is then
determined by the depth percentile of its centroid, providing the
depth-ordered layering that downstream composition requires. This
percentile-based assignment is preferable to a per-pixel depth sort
because it treats each object as an atomic unit, preserving its
identity under subsequent rigid transformation.

\begin{figure}[t]
    \centering
    \resizebox{\linewidth}{!}{%
    \begin{tikzpicture}[
        font=\scriptsize,
        bx/.style={draw=black!55, rounded corners=3pt, thick,
            minimum width=2.0cm, minimum height=0.8cm, align=center},
        ar/.style={->, >=Stealth, thick, draw=black!55},
        adapt/.style={->, >=Stealth, thick, draw=orange!75!black, dashed}
    ]
        \node[bx, fill=yellow!10] (pin) {Painting $P$};
        \node[bx, fill=blue!12, right=0.6cm of pin, minimum width=2.3cm]
            (backbone) {Seg.\ backbone};
        \node[bx, fill=blue!8, below=0.25cm of backbone, minimum width=2.3cm]
            (neck) {Feature pyramid};
        \node[bx, fill=blue!8, below=0.25cm of neck, minimum width=2.3cm]
            (head) {Detection + mask heads};
        \node[bx, fill=gray!8, right=0.6cm of head, minimum width=2.0cm]
            (out) {Masks $\{M_i\}$};

        \node[bx, fill=green!10, right=0.5cm of backbone, minimum width=2.2cm]
            (depth) {Monocular depth};
        \node[bx, fill=gray!8, right=0.5cm of depth, minimum width=1.9cm]
            (dm) {$D_\text{mono}$};

        \node[bx, fill=orange!15, below=0.5cm of head, xshift=-1.2cm,
              minimum width=2.3cm] (adapt)
            {Art-domain\\adaptation};
        \node[bx, fill=purple!10, below=0.5cm of out, xshift=-0.4cm,
              minimum width=2.3cm] (layer)
            {Percentile\\layer assignment};
        \node[bx, fill=gray!8, below=0.25cm of layer, minimum width=2.6cm]
            (final) {$\{M_i\},\,M_\text{bg},\,\{d_i\}$};

        \draw[ar] (pin) -- (backbone);
        \draw[ar] (backbone) -- (neck);
        \draw[ar] (neck) -- (head);
        \draw[ar] (head) -- (out);
        \draw[ar] (pin) -- (depth);
        \draw[ar] (depth) -- (dm);
        \draw[adapt] (backbone.south west) to[bend right=20] (adapt.north);
        \draw[adapt] (adapt) -- (neck);
        \draw[ar] (out) -- (layer);
        \draw[ar] (dm) |- (layer);
        \draw[ar] (layer) -- (final);
    \end{tikzpicture}}
    \caption{\textbf{PaintParse.} A segmentation backbone and a monocular
    depth estimator run in parallel; art-domain adaptation injects
    painting-specific statistics into the backbone, and a percentile-based
    layer assignment fuses masks with depth into a depth-ordered output.}
    \label{fig:paintparse}
\end{figure}

% =====================================================================
\subsection{MonoSplat: Per-Object 3D Lifting and Library Assembly}
\label{sec:method:monosplat}

\paragraph{Design rationale.}
A whole-scene 3D reconstruction would entangle object geometry with
background texture, making independent repositioning impossible. We
therefore lift each foreground object \emph{separately} into a 3D
representation and assemble a \emph{reusable library}, indexed by
semantic label and original painting coordinates. Two design choices
follow from this principle. First, we use a \emph{feed-forward
single-image 3D predictor} rather than per-scene optimization, because
per-object optimization would be prohibitively expensive and would
overfit to a single viewpoint. Second, because feed-forward predictors
operate in a normalized canonical space without metric scale, we
\emph{recover metric dimensions from the PaintParse depth cue}, so that
objects from different paintings can later be composed in a common
metric frame.

\paragraph{Feed-forward lifting.}
For each foreground object $i$, we crop the painting to its bounding
box, $P_i = P \odot M_i$, and feed the crop into a feed-forward
single-image 3D predictor that regresses a set of $N_i$ 3D Gaussians
directly from one image. Each Gaussian is parameterized by a center
$\boldsymbol{\mu}_j \in \mathbb{R}^3$, axis-aligned scales
$\mathbf{s}_j \in \mathbb{R}^3_+$, rotation quaternion
$\mathbf{q}_j \in \mathbb{R}^4$, opacity $\alpha_j \in [0,1]$, and
view-dependent color coefficients. The predictor's generative prior,
trained on diverse 3D assets, generalizes to the missing side and back
of objects that are not visible in the painting---a critical capability
for downstream rotation during composition.

\paragraph{Metric rescaling.}
The predictor outputs in a normalized canonical space. We recover the
metric scale using the depth cue from PaintParse:
\begin{equation}
    \lambda_i \;=\; \frac{d_i \cdot h_i^\text{px}}{f},
    \label{eq:metric_scale}
\end{equation}
where $d_i$ is the median depth of object $i$, $h_i^\text{px}$ is its
pixel height, and $f$ is the focal length in pixels. Each Gaussian
center and scale is then rescaled by $\lambda_i$, producing a metric
3D object. The resulting library
$\{\mathcal{G}_1, \dots, \mathcal{G}_K\}$ is keyed by semantic label
and original painting coordinates, enabling selective retrieval and
cross-painting composition in future work.

\begin{figure}[t]
    \centering
    \resizebox{\linewidth}{!}{%
    \begin{tikzpicture}[
        font=\scriptsize,
        bx/.style={draw=black!55, rounded corners=3pt, thick,
            minimum width=2.1cm, minimum height=0.85cm, align=center},
        ar/.style={->, >=Stealth, thick, draw=black!55}
    ]
        \node[bx, fill=gray!8] (crop) {Crop $P \odot M_i$};
        \node[bx, fill=green!12, right=0.5cm of crop, minimum width=2.7cm,
              minimum height=1.4cm] (lift)
            {\textbf{Feed-forward}\\
             \textbf{single-image 3D}\\
             \footnotesize generative prior};
        \node[bx, fill=gray!8, right=0.5cm of lift, minimum width=2.0cm]
            (raw) {Canonical 3D\\$\widetilde{\mathcal{G}}_i$};
        \node[bx, fill=green!10, below=0.5cm of raw, minimum width=2.4cm]
            (rescale) {Metric rescale\\$\lambda_i = d_i h_i^\text{px}/f$};
        \node[bx, fill=gray!8, below=0.4cm of rescale, minimum width=2.6cm]
            (lib) {Library\\$\{\mathcal{G}_1, \dots, \mathcal{G}_K\}$};
        \node[bx, fill=blue!8, below=0.5cm of lift, minimum width=2.4cm]
            (depth) {Depth cue $d_i$};

        \draw[ar] (crop) -- (lift);
        \draw[ar] (lift) -- (raw);
        \draw[ar] (raw) -- (rescale);
        \draw[ar] (rescale) -- (lib);
        \draw[ar] (depth) -- (rescale);
        \draw[ar] (lift) -- (depth);
    \end{tikzpicture}}
    \caption{\textbf{MonoSplat.} Each foreground crop is lifted into a
    canonical 3D representation by a feed-forward predictor; the
    PaintParse depth cue is used to recover metric scale, and the
    resulting objects are stored in a reusable, retrievable library.}
    \label{fig:monosplat}
\end{figure}

% =====================================================================
\subsection{SketchComp: Line-Art-Guided 3D Scene Composition}
\label{sec:method:sketchcomp}

\paragraph{Design rationale.}
The user's compositional intent---which object goes where---is most
naturally expressed as a sketch, but existing diffusion-based sketch
control treats sketches as appearance or boundary cues rather than
spatial constraints. SketchComp inverts this convention: a sketch
\emph{is} the target, and a 3D object is \emph{geometrically
registered} to it. Two sub-problems must then be solved:
\emph{(i)} finding the rigid transformation that aligns each 3D
object's rendered silhouette with its sketched contour, and
\emph{(ii)} filling the disocclusion gaps that open up when an object
moves away from its original position, while preserving the
background's artistic texture. We solve the first with a
silhouette-overlap objective that combines region IoU with boundary
Chamfer distance, and the second with a depth-conditioned inpainting
module whose geometric conditioning prevents the texture-smear
artifact that naive inpainting produces on paintings.

\paragraph{Contour-based rigid registration.}
Each sketch $\ell_i$ is converted to a binary edge map via Canny
detection, and its external contour $\Gamma_i$ is extracted. For each
3D object $\mathcal{G}_i$, we search for the rigid
transformation---rotation $\theta$, 2D translation $\mathbf{t}$, and
uniform scale $s$---that best aligns the rendered silhouette with
$\Gamma_i$:
\begin{align}
    (\theta_i^*,\, \mathbf{t}_i^*,\, s_i^*) &=
      \arg\min_{\theta,\mathbf{t},s}\;
      \mathcal{L}_\text{cont}(\theta,\mathbf{t},s), \label{eq:contour_reg}\\
    \mathcal{L}_\text{cont} &=
      \underbrace{1 - \text{IoU}\bigl(\mathcal{R}_i,\,\Gamma_i\bigr)}_{\text{region}}
      + \lambda_c \,
      \underbrace{\text{Chamfer}\bigl(\partial\mathcal{R}_i,\,\Gamma_i\bigr)}_{\text{boundary}},
      \nonumber
\end{align}
where $\mathcal{R}_i = \mathcal{R}(\mathcal{G}_i;\theta,\mathbf{t},s)$
is the silhouette rendered from the canonical camera and $\lambda_c$
balances the two terms. The IoU term rewards global region agreement,
while the Chamfer term penalizes boundary misalignment---their
combination avoids the degenerate solutions that either term alone
admits (a small silhouette inside the contour maximizes IoU locally;
boundary-only matching can produce topologically wrong regions).
Optimization proceeds via gradient descent with iterative
re-rendering, so that each step's gradient is computed against the
current geometry rather than a stale approximation.

\paragraph{Scene assembly and disocclusion inpainting.}
Each registered object is placed at its PaintParse-derived depth
$d_i$ and objects are rendered in depth-sorted order against a
background layer reconstructed from the RGB-D map. Object
displacement creates \emph{disocclusion gaps}---background pixels
newly visible after objects move---which we detect as
\begin{equation}
    \Omega_\text{gap} \;=\;
    \bigl\{\,p \in M_\text{bg}^\text{orig} \;\big|\;
        \text{no 3D surface projects to } p\,\bigr\}.
    \label{eq:gap}
\end{equation}
Filling $\Omega_\text{gap}$ while preserving artistic texture coherence
is critical. We employ a depth-conditioned inpainting diffusion
pipeline: a masked composite (with $\Omega_\text{gap}$ zeroed) is fed
into a latent inpainting diffusion model with a depth-conditioned
control branch conditioned on the composited depth map. The depth
conditioning is the key design choice: it constrains the inpainted
content to follow the background's geometric structure, preventing
the texture-smear artifact that arises when an unconstrained
inpainter fills a painting gap with semantically plausible but
geometrically inconsistent content. A dilation radius expands the
surrounding context to supply texture reference, ensuring the filled
region blends with adjacent brushwork. The resulting scene
$\mathcal{S}^*$ and its depth map $\mathcal{D}^*$ are forwarded to
ArtiFlow.

\begin{figure}[t]
    \centering
    \resizebox{\linewidth}{!}{%
    \begin{tikzpicture}[
        font=\scriptsize,
        bx/.style={draw=black!55, rounded corners=3pt, thick,
            minimum width=2.1cm, minimum height=0.85cm, align=center},
        block/.style={draw=black!55, rounded corners=3pt, thick,
            fill=orange!14, minimum width=2.4cm, minimum height=1.2cm,
            align=center},
        inpblock/.style={draw=black!55, rounded corners=3pt, thick,
            fill=purple!14, minimum width=2.4cm, minimum height=1.2cm,
            align=center},
        ar/.style={->, >=Stealth, thick, draw=black!55}
    ]
        \node[bx, fill=yellow!10] (la) {Line-art $L$};
        \node[bx, fill=green!10, below=0.3cm of la, minimum width=2.2cm]
            (lib) {3D library $\{\mathcal{G}_i\}$};
        \node[block, right=0.6cm of la] (canny)
            {\textbf{Contour extraction}\\
             \footnotesize Canny + external contour $\Gamma_i$};
        \node[block, right=0.5cm of canny] (reg)
            {\textbf{Rigid registration}\\
             \footnotesize IoU + Chamfer\\
             \footnotesize (Eq.~\ref{eq:contour_reg})};
        \node[bx, fill=gray!8, right=0.5cm of reg, minimum width=2.1cm]
            (placed) {Placed $\{\mathcal{G}'_i\}$};
        \node[block, below=0.55cm of reg] (assemble)
            {\textbf{Depth-sorted assembly}\\
             \footnotesize render $\mathcal{S}$, $\mathcal{D}$};
        \node[bx, fill=gray!8, right=0.5cm of assemble, minimum width=2.0cm]
            (gap) {Gap $\Omega_\text{gap}$\\(Eq.~\ref{eq:gap})};
        \node[inpblock, below=0.55cm of assemble] (inp)
            {\textbf{Depth-conditioned inpainting}\\
             \footnotesize inpaint diffusion + depth control};
        \node[bx, fill=gray!8, right=0.5cm of inp, minimum width=2.1cm]
            (fin) {$\mathcal{S}^*,\;\mathcal{D}^*$};

        \draw[ar] (la) -- (canny);
        \draw[ar] (canny) -- (reg);
        \draw[ar] (lib) -| (reg);
        \draw[ar] (reg) -- (placed);
        \draw[ar] (placed) -- (assemble);
        \draw[ar] (assemble) -- (gap);
        \draw[ar] (gap) -- (inp);
        \draw[ar] (inp) -- (fin);
        \draw[ar] (assemble) |- (fin);
    \end{tikzpicture}}
    \caption{\textbf{SketchComp.} Hand-drawn contours are extracted and
    matched to rendered 3D silhouettes via a rigid registration that
    jointly optimizes region IoU and boundary Chamfer distance. After
    depth-sorted scene assembly, disocclusion gaps are detected and
    filled by a depth-conditioned inpainting diffusion module that
    respects the background's geometric structure.}
    \label{fig:sketchcomp}
\end{figure}

% =====================================================================
\subsection{ArtiFlow: Geometry-Grounded Video Synthesis}
\label{sec:method:artiflow}

\paragraph{Design rationale.}
Pure image-to-video diffusion applied to a painting produces
plausible per-frame motion but suffers two failure modes: geometric
drift (objects wander off their intended trajectories) and boundary
melting (brushstroke edges blur across frames). ArtiFlow addresses
both by \emph{grounding the diffusion in explicit 3D geometry}: the
composited scene is rendered into a multi-view keyframe stream that
provides three control signals---RGB appearance, depth geometry, and
edge boundary---which are injected into a video diffusion backbone as
dual control branches. The depth branch anchors each frame to the 3D
structure of $\mathcal{S}^*$, suppressing geometric drift; the edge
branch preserves brushstroke-level boundary detail, suppressing
boundary melting. The combined conditioning converts the diffusion
backbone from a free-form motion generator into a \emph{geometry-anchored
refiner} whose temporal freedom is restricted to the residual motion
that the 3D rendering does not already specify.

\paragraph{Multi-view keyframe rendering.}
Given $\mathcal{S}^*$, we define a smooth camera trajectory through
$T$ poses $\{\mathbf{T}_t\}_{t=1}^{T}$, generated by cubic B-spline
interpolation between a canonical front-facing start pose and a
user-specified or auto-generated end pose. At each step $t$,
differentiable rasterization produces an RGB frame
$I_t = \mathcal{R}_\text{rgb}(\mathcal{S}^*,\mathbf{T}_t)$, a depth
map $D_t = \mathcal{R}_\text{depth}(\mathcal{S}^*,\mathbf{T}_t)$, and
an edge map $E_t = \text{Edge}(I_t)$. The keyframe count is set to
$T = 8N+1$ to align with the temporal chunk constraints of modern
video diffusion backbones.

\paragraph{Dual-conditioned video diffusion.}
Two parallel control branches inject geometric and boundary cues
into the video diffusion backbone. The depth branch, conditioned on
$\{D_t\}$ at strength $s_d$, anchors each frame to the 3D structure
of $\mathcal{S}^*$, preventing the geometric drift that plagues pure
I2V models. The edge branch, conditioned on $\{E_t\}$ at strength
$s_e$, preserves brushstroke-level boundary detail, suppressing the
boundary-melting artifact common when diffusion is applied to
paintings. The combined conditioning signal per frame is
\begin{equation}
    \mathbf{c}_t \;=\; s_d \, \mathbf{c}_t^\text{depth}
                   \;+\; s_e \, \mathbf{c}_t^\text{edge},
    \label{eq:dual_control}
\end{equation}
injected at the corresponding residual connections of the backbone.
Temporal coherence is further reinforced by \emph{first-frame latent
anchoring}: subsequent frames' initial noise is set to the first-frame
latent plus a position-encoded Gaussian perturbation
$\boldsymbol{\epsilon}_t \sim \mathcal{N}(0, \sigma_t^2 \mathbf{I})$
with $\sigma_t$ growing linearly with $|t-1|$, so that motion freedom
increases smoothly away from the anchored first frame.

\paragraph{Per-object motion verification.}
After generation, we verify motion fidelity by computing the IoU
between rendered object masks and segmented masks in each generated
frame. If the average IoU falls below a threshold $\tau_\text{mo}$,
the pipeline triggers a \emph{re-composition} step (dashed arrow in
Figure~\ref{fig:overview}) that adjusts control strengths or
regenerates the affected segment with stronger depth conditioning,
closing the loop between composition and synthesis.

\begin{figure}[t]
    \centering
    \resizebox{\linewidth}{!}{%
    \begin{tikzpicture}[
        font=\scriptsize,
        bx/.style={draw=black!55, rounded corners=3pt, thick,
            minimum width=2.0cm, minimum height=0.8cm, align=center},
        block/.style={draw=black!55, rounded corners=3pt, thick,
            fill=red!12, minimum width=2.4cm, minimum height=1.2cm,
            align=center},
        ar/.style={->, >=Stealth, thick, draw=black!55}
    ]
        \node[bx, fill=purple!10] (sc) {$\mathcal{S}^*,\,\mathcal{D}^*$};
        \node[block, right=0.5cm of sc, minimum height=1.4cm] (rend)
            {\textbf{Multi-view rendering}\\
             \footnotesize B-spline trajectory\\
             \footnotesize $\{I_t, D_t, E_t\}$};
        \node[bx, fill=gray!8, right=0.45cm of rend, minimum width=1.6cm]
            (rgb) {RGB $I_t$};
        \node[bx, fill=blue!10, below=0.2cm of rgb, minimum width=1.6cm]
            (dep) {Depth $D_t$};
        \node[bx, fill=orange!12, below=0.2cm of dep, minimum width=1.6cm]
            (edg) {Edge $E_t$};

        \node[block, below=0.6cm of rend, xshift=0.4cm,
              minimum width=3.0cm] (cn)
            {\textbf{Dual control branches}\\
             \footnotesize depth + edge injection (Eq.~\ref{eq:dual_control})};
        \node[block, right=0.4cm of cn, minimum width=2.6cm,
              minimum height=1.4cm] (vd)
            {\textbf{Video diffusion backbone}\\
             \footnotesize temporal attention\\
             \footnotesize first-frame anchor};
        \node[bx, fill=red!18, below=0.45cm of vd, minimum width=2.4cm]
            (out) {Output video $V$};

        \draw[ar] (sc) -- (rend);
        \draw[ar] (rend) -- (rgb);
        \draw[ar] (rend) -- (dep);
        \draw[ar] (rend) -- (edg);
        \draw[ar] (dep) -- (cn);
        \draw[ar] (edg) -- (cn);
        \draw[ar] (cn) -- (vd);
        \draw[ar] (rgb) |- (vd);
        \draw[ar] (vd) -- (out);
    \end{tikzpicture}}
    \caption{\textbf{ArtiFlow.} The composited 3D scene is rendered into
    a multi-view keyframe stream providing RGB, depth, and edge signals.
    A dual-branch control injection (depth + edge) anchors the video
    diffusion backbone to the 3D geometry, with first-frame latent
    anchoring for temporal coherence.}
    \label{fig:artiflow}
\end{figure}

% =====================================================================
\subsection{Pipeline Algorithm and Implementation}
\label{sec:method:algo}

Algorithm~\ref{alg:canvas3d} summarizes the full Canvas3D pipeline. The
modular design enables \emph{fault-tolerant resumption}: each stage
saves its outputs as checkpoint files, so any specialist can be
upgraded or replaced without re-running downstream stages---for
example, swapping in a stronger 3D predictor or a different video
backbone. Peak memory consumption across all four stages is kept
within a single 24\,GB budget through serial model loading with
explicit memory release between stages, which is feasible precisely
because the handshake representations between stages are small,
serializable tensors rather than shared activations.

\begin{algorithm}[t]
\caption{Canvas3D Pipeline}
\label{alg:canvas3d}
{\footnotesize
\begin{algorithmic}[1]
\REQUIRE Painting $P$; line-art sketches $L=\{\ell_i\}_{i=1}^{K}$;
         thresholds $\lambda_c$, $s_d$, $s_e$, $\tau_\text{mo}$
\ENSURE Animated video $V$ with per-object motion
\medskip
\STATE \textbf{// Stage 1 --- PaintParse: perceptual grouping + layers}
\STATE $\{M_i\}_{i=1}^{K} \gets \text{SegAdapt}(P)$
    \hfill $\triangleright$ art-domain-adapted segmentation
\STATE $M_\text{bg} \gets 1 - \max_i M_i$
    \hfill $\triangleright$ Eq.~\ref{eq:bg_mask}
\STATE $D_\text{mono} \gets \text{DepthEstimator}(P)$
\STATE \textbf{for} $i = 1$ \textbf{to} $K$ \textbf{do}
\STATE \quad $\Omega_i \gets \{p \mid M_i(p){=}1\}$;
    $d_i \gets \text{Median}\{D_\text{mono}(p) \mid p \in \Omega_i\}$
\STATE \textbf{end for}
\medskip
\STATE \textbf{// Stage 2 --- MonoSplat: per-object 3D lifting}
\STATE \textbf{for} $i = 1$ \textbf{to} $K$ \textbf{do}
\STATE \quad $\widetilde{\mathcal{G}}_i \gets
    \text{FeedForward3D}(P \odot M_i)$
    \hfill $\triangleright$ canonical-space 3D
\STATE \quad $\lambda_i \gets d_i \cdot h_i^\text{px}/f$
    \hfill $\triangleright$ Eq.~\ref{eq:metric_scale}
\STATE \quad $\mathcal{G}_i \gets \lambda_i \cdot \widetilde{\mathcal{G}}_i$
    \hfill $\triangleright$ metric rescaling
\STATE \quad $\text{Lib}[i] \gets (\mathcal{G}_i,\, d_i,\, M_i)$
\STATE \textbf{end for}
\medskip
\STATE \textbf{// Stage 3 --- SketchComp: composition + inpainting}
\STATE \textbf{for} $i = 1$ \textbf{to} $K$ \textbf{do}
\STATE \quad $\Gamma_i \gets \text{ExternalContour}(\text{Canny}(\ell_i))$
\STATE \quad $(\theta_i^*,\mathbf{t}_i^*,s_i^*) \gets
    \arg\min \mathcal{L}_\text{cont}(\mathcal{G}_i,\Gamma_i)$
    \hfill $\triangleright$ Eq.~\ref{eq:contour_reg}
\STATE \quad $\mathcal{G}'_i \gets
    \text{Transform}(\mathcal{G}_i,\theta_i^*,\mathbf{t}_i^*,s_i^*)$
\STATE \textbf{end for}
\STATE $\mathcal{S},\mathcal{D} \gets
    \text{DepthSort}(\{\mathcal{G}'_i\},\,M_\text{bg},\,\{d_i\})$
\STATE $\Omega_\text{gap} \gets
    \text{DetectGap}(\mathcal{S},\,M_\text{bg}^\text{orig})$
    \hfill $\triangleright$ Eq.~\ref{eq:gap}
\STATE $\mathcal{S}^*,\mathcal{D}^* \gets
    \text{DepthInpaint}(\mathcal{S},\,\mathcal{D},\,\Omega_\text{gap})$
\medskip
\STATE \textbf{// Stage 4 --- ArtiFlow: geometry-grounded video synthesis}
\STATE $\{\mathbf{T}_t\}_{t=1}^{T} \gets \text{BSplineTrajectory}(T)$
\STATE \textbf{for} $t = 1$ \textbf{to} $T$ \textbf{do}
\STATE \quad $I_t \gets \mathcal{R}_\text{rgb}(\mathcal{S}^*,\mathbf{T}_t)$;
\STATE \quad $D_t \gets \mathcal{R}_\text{depth}(\mathcal{S}^*,\mathbf{T}_t)$;
\STATE \quad $E_t \gets \text{Edge}(I_t)$
\STATE \textbf{end for}
\STATE $V \gets \text{VideoDiffusion}\bigl(\{I_t\},\,s_d\!\cdot\!D_t
    + s_e\!\cdot\!E_t\bigr)$
    \hfill $\triangleright$ Eq.~\ref{eq:dual_control}
\STATE \textbf{if} $\text{MotionIoU}(V) < \tau_\text{mo}$ \textbf{then}
    \quad \textbf{go to} Stage 3
    \hfill $\triangleright$ re-composition loop
\STATE \textbf{return} $V$
\end{algorithmic}}
\end{algorithm}
